\clearpage
\VerdierPhysicalPageAnchor{0198}%
\begingroup\small\setlength{\parskip}{0pt}\setlength{\emergencystretch}{1.5em}
\setlength{\abovedisplayskip}{4pt plus 2pt minus 2pt}%
\setlength{\belowdisplayskip}{4pt plus 2pt minus 2pt}%

\VerdierSemanticAnchor{ch3d-2-3-1-proof-b-conclusion}
and hence $v$ is an epimorphism, which proves property $(E_3)$ and completes
the proof of the proposition.

\VerdierSemanticAnchor{ch3d-2-3-2}
\medskip\noindent\textbf{Proposition 2.3.2.}
\textit{Let $M$ be a set of objects of $\mathcal A$ satisfying properties
$(E_1)$, $(E_2)$, and $(E_3)$. Let
$Y\in\operatorname{Ob}\operatorname{comp}(\mathcal A)$ be a complex such
that, for every integer $n$, $H^n(Y)\in M$, and suppose that there is an
integer $n_0$ such that $H^n(Y)=0$ for $n\geq n_0$. Then there is a complex
$X\in\operatorname{Ob}\operatorname{comp}^{-}(\mathcal A)$ whose components
belong to $M$, and a quasi-isomorphism}
\[
\dot u\colon X\longrightarrow Y.
\]

\VerdierSemanticAnchor{ch3d-2-3-2-proof}
First, by (1.2.1), there is a complex
$Y'\in\operatorname{Ob}\operatorname{comp}^{-}(\mathcal A)$ and a
quasi-isomorphism $Y'\longrightarrow Y$. We may therefore reduce to the case
where $Y$ is an object of $\operatorname{comp}^{-}(\mathcal A)$. We shall
construct inductively on $n$ a sequence of morphisms

\VerdierSemanticAnchor{ch3d-2-3-2-induction}
\[
u_n\colon X_n\longrightarrow X_{n-1}
\]
having the following properties:

\VerdierSemanticAnchor{ch3d-2-3-2-induction-properties}
\noindent 1) For $n$ sufficiently small, $X_n=Y$ and
$u_n=\operatorname{id}$.

\noindent 2) The components of degree $i\geq -n$ of the complex $X_n$ belong
to $M$. The kernels and images of the differentials of degree $i\geq -n$ of
$X_n$ belong to $M^\wedge$ (2.1.1).

\noindent 3) The morphism $u_n$ is a quasi-isomorphism. It is isomorphic to
the morphism $p((\dot u_n)^{-n},-n)$ (2.1.1.3).

\VerdierSemanticAnchor{ch3d-2-3-2-limits}
Once this has been done, it suffices to set $X=\varprojlim_n X_n$ and take
as $\dot u$ the canonical morphism
\[
\dot u\colon X=\varprojlim_n X_n\longrightarrow\varinjlim_n X_n=Y.
\]

\VerdierSemanticAnchor{ch3d-2-3-2-induction-step}
For $n$ sufficiently small, set $X_n=Y$ and
$\dot u_n=\operatorname{id}_Y$. Suppose that, for every integer $n<q$,
complexes $X_n$ and morphisms $\dot u_n$ having properties 1), 2), and 3) above
have been constructed. The complex $X_{q-1}$ has the form

\VerdierSemanticAnchor{ch3d-2-3-2-xqminus1}
\[
X_{q-1}=\cdots\longrightarrow Y^{-q-2}\longrightarrow Y^{-q-1}
\longrightarrow Z^{-q}\longrightarrow L^{-q+1}\longrightarrow L^{-q+2}
\longrightarrow\cdots,
\]
where the $Y^i$ are the components of the complex $Y$ and the $L^i$ belong
to $M$. We then have an epimorphism

\VerdierSemanticAnchor{ch3d-2-3-2-quotient}
\[
Z^{-q}\longrightarrow
Z^{-q}/\operatorname{Im}(Y^{-q-1}\longrightarrow Z^{-q})
\longrightarrow0
\]
\endgroup
